% ============================================================
%  Ângulos: Relações Angulares — Apostila Premium | PratiqueJá
%  Engine: XeLaTeX
% ============================================================
\documentclass[12pt]{article}
\usepackage{../../pratiqueja}
\usepackage{tikz}
\usetikzlibrary{calc}

% \hlm — destaque de resultado em modo-math (cor sem bold)
\newcommand{\hlm}[1]{{\color{darkblue}#1}}

\setsubject{Relações Angulares}

\begin{document}
\pagenumbering{arabic}
\setstretch{1.2}
\color{bodytext}

\heroheader{Relações entre Ângulos}%
           {Opostos, Complementares, Suplementares e Paralelas}%
           {Matemática · Intermediário}

\setlength{\columnsep}{18pt}
\setlength{\columnseprule}{0.5pt}
\def\columnseprulecolor{\color{exborder}}

\begin{multicols}{2}

%─────────────────────────────────────────────────────────────
\TW{Def}{badgepurple}{Opostos pelo Vértice}{Também chamados de ângulos verticais}

\regra{Dois ângulos \textbf{opostos pelo vértice} têm sempre a mesma
medida — são \textbf{congruentes}.}

\begin{center}\vspace{2pt}
\begin{tikzpicture}[scale=0.8]
\draw[line width=1pt, darkblue] (60:2.1)--(240:2.1);
\draw[line width=1pt, darkblue] (120:2.1)--(300:2.1);
\draw[vered, line width=1.4pt] (60:0.7) arc (60:120:0.7);
\draw[vered, line width=1.4pt] (240:0.7) arc (240:300:0.7);
\draw[vegreen, line width=1.4pt] (-60:0.7) arc (-60:60:0.7);
\draw[vegreen, line width=1.4pt] (120:0.7) arc (120:240:0.7);
\node[vered, font=\small\bfseries] at (90:1.05){$\alpha$};
\node[vered, font=\small\bfseries] at (270:1.05){$\alpha$};
\node[vegreen, font=\small\bfseries] at (0:1.2){$\beta$};
\node[vegreen, font=\small\bfseries] at (180:1.2){$\beta$};
\end{tikzpicture}
\end{center}

\obs{Os pares \textbf{não adjacentes} são opostos pelo vértice ($\alpha$
e $\alpha$, $\beta$ e $\beta$). Já $\alpha$ e $\beta$ são adjacentes e
\textbf{suplementares}: $\alpha + \beta = 180°$.}

\SH{Exemplo — $\alpha = 60°$}
\exemp{%
  $x = \alpha = \hlm{60°}$ {\footnotesize(oposto → congruente)}\\[2pt]
  $\beta = 180° - 60° = \hlm{120°}$ {\footnotesize(adjacente → suplementar)}%
}

%─────────────────────────────────────────────────────────────
\TW{90°}{badgegreen}{Ângulos Complementares}{Soma igual a $90°$}

\regra{Dois ângulos são \textbf{complementares} quando somam $90°$.
O complemento de $\alpha$ é $90° - \alpha$.}

\obs{Não precisam ser adjacentes. Aparecem nos \textbf{triângulos
retângulos}, onde os ângulos agudos são sempre complementares.}

\exemp{%
  $30° + 60° = \hlm{90°}$ \ok \qquad $45° + 45° = \hlm{90°}$ \ok\\[2pt]
  $50° + 60° = 110° \neq 90°$ \nok%
}

\SH{Encontrando o complemento}
\exemp{%
  $\alpha = 37° \Rightarrow 90° - 37° = \hlm{53°}$\\[2pt]
  $\alpha = 72° \Rightarrow 90° - 72° = \hlm{18°}$%
}

%─────────────────────────────────────────────────────────────
\TW{180°}{vered}{Ângulos Suplementares}{Soma igual a $180°$}

\regra{Dois ângulos são \textbf{suplementares} quando somam $180°$.
O suplemento de $\alpha$ é $180° - \alpha$.}

\obs{Suplementares adjacentes formam uma \textbf{linha reta} — cada
ângulo é suplementar a seus vizinhos adjacentes.}

\exemp{%
  $60° + 120° = \hlm{180°}$ \ok \qquad $90° + 90° = \hlm{180°}$ \ok\\[2pt]
  $70° + 80° = 150° \neq 180°$ \nok%
}

\SH{Encontrando o suplemento}
\exemp{%
  $\alpha = 30° \Rightarrow 180° - 30° = \hlm{150°}$\\[2pt]
  $\alpha = 110° \Rightarrow 180° - 110° = \hlm{70°}$%
}

%─────────────────────────────────────────────────────────────
\TW{$\parallel$}{darkblue}{Paralelas e Transversal}{Correspondentes, alternos e colaterais}

\regra{Uma transversal cortando duas paralelas forma \textbf{8 ângulos}:
\textbf{correspondentes} e \textbf{alternos} são iguais; \textbf{colaterais}
são suplementares ($180°$).}

\begin{center}\vspace{2pt}
\begin{tikzpicture}[scale=0.72]
\coordinate (P1) at (0.63,1.4);
\coordinate (P2) at (-0.63,-1.4);
\draw[line width=1pt, darkblue] (-2.8,1.4)--(2.8,1.4) node[right,font=\footnotesize\bfseries]{$r$};
\draw[line width=1pt, darkblue] (-2.8,-1.4)--(2.8,-1.4) node[right,font=\footnotesize\bfseries]{$s$};
\draw[line width=1pt, graycolor] (0.945,2.1)--(-0.945,-2.1) node[below,font=\footnotesize\bfseries]{$t$};
% P1
\draw[vered, line width=1.1pt] ([shift=(0:0.45)]P1) arc (0:66:0.45);
\draw[vegreen, line width=1.1pt] ([shift=(66:0.4)]P1) arc (66:180:0.4);
\node[vered, font=\footnotesize\bfseries] at ($(P1)+(33:0.74)$){$\alpha$};
\node[vegreen, font=\footnotesize\bfseries] at ($(P1)+(123:0.66)$){$\beta$};
% P2
\draw[vered, line width=1.1pt] ([shift=(0:0.45)]P2) arc (0:66:0.45);
\draw[vegreen, line width=1.1pt] ([shift=(66:0.4)]P2) arc (66:180:0.4);
\node[vered, font=\footnotesize\bfseries] at ($(P2)+(33:0.74)$){$\alpha$};
\node[vegreen, font=\footnotesize\bfseries] at ($(P2)+(123:0.66)$){$\beta$};
\end{tikzpicture}
\end{center}

\SH{Tipos de ângulos formados}
\exemp{%
  \textbf{Correspondentes:} mesmo lado da transversal, um em cada
  paralela → \hlm{iguais}\\[2pt]
  \textbf{Alternos} (internos/externos): lados \emph{opostos} da
  transversal → \hlm{iguais}\\[2pt]
  \textbf{Colaterais} (co-internos/externos): \emph{mesmo lado} da
  transversal → \hlm{suplementares}%
}

\nota{\textbf{Erro comum:} alternos internos e colaterais co-internos
ficam ambos \emph{entre} as paralelas, mas alternos são \textbf{iguais}
(lados opostos) e colaterais são \textbf{suplementares} (mesmo lado).}

\SH{Exemplos resolvidos}
\exemp{%
  \textbf{1) Correspondente} — $\alpha = 50°$ em $r$:\\
  correspondente em $s = \hlm{50°}$; adjacente $\beta = 180°-50° = \hlm{130°}$\\[3pt]
  \textbf{2) Alterno interno} — $\alpha = 70°$ em $r$:\\
  alterno em $s = \hlm{70°}$ {\footnotesize(iguais — lados opostos)}\\[3pt]
  \textbf{3) Colateral co-interno} — $\alpha = 30°$ em $r$:\\
  colateral em $s = 180°-30° = \hlm{150°}$ {\footnotesize(suplementares)}%
}

\obs{\textbf{Teorema recíproco:} se dois ângulos correspondentes (ou
alternos) são iguais, então as retas cortadas são \textbf{paralelas} —
usado para provar paralelismo.}

\end{multicols}

%─────────────────────────────────────────────────────────────
\vspace{4pt}
\TW{$\equiv$}{badgegreen}{Tabela Resumo}{Todas as relações angulares}

\vspace{6pt}
\begin{center}
\setlength{\tabcolsep}{8pt}\renewcommand{\arraystretch}{1.3}\small
\begin{tabular}{llcl}
\rowcolor{darkblue}
\textcolor{white}{\textbf{Tipo}} & \textcolor{white}{\textbf{Contexto}} &
\textcolor{white}{\textbf{Fórmula}} & \textcolor{white}{\textbf{Exemplo}}\\
Opostos pelo vértice   & Interseção de duas retas   & $\alpha_1 = \alpha_2$      & Se $\alpha=50°$, o oposto é $50°$\\
\rowcolor{exbg}
Complementares         & Par cuja soma é $90°$      & $\alpha + \beta = 90°$     & $30° + 60° = 90°$\\
Suplementares          & Par cuja soma é $180°$     & $\alpha + \beta = 180°$    & $60° + 120° = 180°$\\
\rowcolor{exbg}
Correspondentes        & Paralelas + transversal    & $\alpha_r = \alpha_s$      & Se $\alpha=70°$, corresp.\ é $70°$\\
Alternos internos      & Paralelas + transversal    & $\alpha_r = \alpha_s$      & Se $\alpha=70°$, alterno é $70°$\\
\rowcolor{exbg}
Alternos externos      & Paralelas + transversal    & $\alpha_r = \alpha_s$      & Se $\alpha=70°$, alt.\ ext.\ é $70°$\\
Colaterais co-internos & Paralelas + transversal    & $\alpha + \beta = 180°$    & Se $\alpha=70°$, colateral é $110°$\\
\rowcolor{exbg}
Colaterais co-externos & Paralelas + transversal    & $\alpha + \beta = 180°$    & Se $\alpha=70°$, col.\ ext.\ é $110°$\\
\end{tabular}
\end{center}

\end{document}
